\section{Elementary Row Operations Revisited}

Let $A \in \M_{n \times p}(K)$. We have previously learned about elementary row operations, which are utilized as an integral part of the Gaussian elimination algorithm. The main point of this section is to demonstrate that all of these operations can be expressed through matrix multiplications.
\textbf{Notation:} Let $n \in \mathbb{Z}_{\ge 1}$, and let $1 \le i \le n$, $1 \le j \le n$. We denote by $E_{ij} \in \M_{n \times n}(K)$ the matrix whose entries are all zero, except for the entry in the $i$-th row and $j$-th column, which is equal to $1$. 

We will now define three types of elementary matrices in $\M_{n \times n}(K)$.

\begin{definition}[Elementary Matrices]
% def:17.e.0
\label{def:elementary_matrices}
We define the following three types of elementary matrices:
\begin{enumerate}
    \item \textbf{Type I:} Let $1 \le i \le n$, $1 \le j \le n$ with $i \neq j$, and let $\alpha \in K$. We define $Q_{ij}(\alpha)$ to be the matrix obtained from $I_n$ by applying the row operation $R_i + \alpha R_j \longrightarrow R_i$ (i.e., replacing the $i$-th row by the sum of the $i$-th row and $\alpha$ times the $j$-th row). Explicitly, we have:
    \begin{equation*}
        Q_{ij}(\alpha) = I_n + \alpha \cdot E_{ij}
    \end{equation*}
    
    \item \textbf{Type II:} Let $1 \le i \le n$, $1 \le j \le n$ with $i \neq j$. We define $P_{ij}$ to be the matrix obtained from $I_n$ by permuting (exchanging) the $i$-th row and the $j$-th row, denoted by $R_i \leftrightarrow R_j$.
    
    \item \textbf{Type III:} Let $1 \le i \le n$, and let $0 \neq \alpha \in K$. We denote by $S_i(\alpha)$ the matrix obtained from $I_n$ by applying the row operation $\alpha R_i \longrightarrow R_i$ (i.e., multiplying the $i$-th row by the non-zero scalar $\alpha$).
\end{enumerate}
\end{definition}

\begin{exercise}
Show that the Type II elementary matrix can be written explicitly as $P_{ij} = I_n - E_{ii} - E_{jj} + E_{ij} + E_{ji}$.
\end{exercise}

\begin{example}
Consider the space $\M_{4 \times 4}(K)$. Examples of the three types of elementary matrices are:
\begin{align*}
    Q_{1,3}(\alpha) &= \begin{pmatrix} 1 & 0 & \alpha & 0 \\ 0 & 1 & 0 & 0 \\ 0 & 0 & 1 & 0 \\ 0 & 0 & 0 & 1 \end{pmatrix} \\
    P_{3,4} &= \begin{pmatrix} 1 & 0 & 0 & 0 \\ 0 & 1 & 0 & 0 \\ 0 & 0 & 0 & 1 \\ 0 & 0 & 1 & 0 \end{pmatrix} \\
    S_3(\alpha) &= \begin{pmatrix} 1 & 0 & 0 & 0 \\ 0 & 1 & 0 & 0 \\ 0 & 0 & \alpha & 0 \\ 0 & 0 & 0 & 1 \end{pmatrix}
\end{align*}
\end{example}

\begin{lemma}[Matrix Multiplication and Row Operations]
% lemma:17.e.1
\label{lemma:row_operations_matrix_multiplication}
Let $A \in \M_{n \times p}(K)$. Multiplying $A$ from the left by an elementary matrix performs the corresponding elementary row operation on $A$:
\begin{enumerate}
    \item Multiplying $A$ from the left with $Q_{ij}(\alpha)$ results in the row operation $R_i + \alpha R_j \longrightarrow R_i$:
    \begin{equation*}
        A \xrightarrow{R_i + \alpha R_j \longrightarrow R_i} Q_{ij}(\alpha) \cdot A
    \end{equation*}
    \item Multiplying $A$ from the left with $P_{ij}$ results in the row operation $R_i \leftrightarrow R_j$:
    \begin{equation*}
        A \xrightarrow{R_i \leftrightarrow R_j} P_{ij} \cdot A
    \end{equation*}
    \item Multiplying $A$ from the left with $S_i(\alpha)$ results in the row operation $\alpha R_i \longrightarrow R_i$:
    \begin{equation*}
        A \xrightarrow{\alpha R_i \longrightarrow R_i} S_i(\alpha) \cdot A
    \end{equation*}
\end{enumerate}
\end{lemma}

\begin{importantremark}[Column Operations]
The same principle holds true for elementary column operations, but via right-multiplication:
\begin{enumerate}
    \item[\textbf{(a')}] Multiplying $A$ from the right with $Q_{ij}(\alpha)$ results in the column operation $C_j + \alpha C_i \longrightarrow C_j$:
    \begin{equation*}
        A \xrightarrow{C_j + \alpha C_i \longrightarrow C_j} A \cdot Q_{ij}(\alpha)
    \end{equation*}
    \item[\textbf{(b')}] Multiplying $A$ from the right with $P_{ij}$ results in the column operation $C_i \leftrightarrow C_j$:
    \begin{equation*}
        A \xrightarrow{C_i \leftrightarrow C_j} A \cdot P_{ij}
    \end{equation*}
    \item[\textbf{(c')}] Multiplying $A$ from the right with $S_i(\alpha)$ results in the column operation $\alpha C_i \longrightarrow C_i$:
    \begin{equation*}
        A \xrightarrow{\alpha C_i \longrightarrow C_i} A \cdot S_i(\alpha)
    \end{equation*}
\end{enumerate}
\end{importantremark}

\begin{proof}
Statements (a), (b), and (c) can be proven directly by calculating the matrix multiplications (you are encouraged to check this!). 

Statements (a'), (b'), and (c') can be deduced from the row operation versions by passing to the transposed matrix. For example, to prove (a'):
\begin{equation*}
    (A \cdot Q_{ij}(\alpha))^T = Q_{ij}(\alpha)^T \cdot A^T = Q_{ji}(\alpha) \cdot A^T
\end{equation*}
By part (a), this is exactly the matrix obtained from $A^T$ after applying to it the row operation $R_j + \alpha R_i \longrightarrow R_j$. But taking the transpose again, this is equal to:
\begin{equation*}
    \left( \text{matrix obtained from } A \text{ after applying } C_j + \alpha C_i \longrightarrow C_j \right)^T
\end{equation*}
Taking the transpose of both sides, it follows that $A \cdot Q_{ij}(\alpha)$ is the matrix obtained from $A$ after the column operation $C_j + \alpha C_i \longrightarrow C_j$. The proofs for (b') and (c') follow a similar logic.
\end{proof}

\begin{lemma}[Inverses of Elementary Matrices]
% lemma:17.e.2
\label{lemma:elementary_matrix_inverses}
Every elementary matrix is invertible, and the inverse of each is an elementary matrix of the exact same type:
\begin{equation*}
    Q_{ij}(\alpha)^{-1} = Q_{ij}(-\alpha), \quad P_{ij}^{-1} = P_{ij}, \quad \text{and} \quad S_i(\alpha)^{-1} = S_i(\alpha^{-1})
\end{equation*}
\end{lemma}

\begin{proof}
\textit{Exercise.} (\textit{Hint:} Note that $P_{ij} \cdot P_{ij}$ is the matrix obtained from $P_{ij}$ by exchanging rows $i$ and $j$, which brings us immediately back to $I_n$. Apply similar logical steps for the other types.)
\end{proof}

\begin{theorem}[Decomposition into Elementary Matrices]
% thm:17.e.3
\label{thm:elementary_matrix_decomposition}
For every invertible matrix $A \in GL_n(K)$, there exists an integer $k \ge 1$ and a sequence of elementary matrices $T_1, \dots, T_k$ such that:
\begin{equation*}
    T_k \cdot T_{k-1} \dots T_1 \cdot A = I_n
\end{equation*}
Furthermore, we have $A = T_1^{-1} \dots T_k^{-1}$ and $A^{-1} = T_k \dots T_1$. In particular, any invertible matrix $A$ can be written as a finite product of elementary matrices.
\end{theorem}

\begin{proof}
We already know that every matrix $A \in \M_{n \times n}(K)$ can be brought to a row-reduced echelon matrix, which we will call $A'$, using the Gaussian elimination algorithm. 

In our case, $A$ is invertible, and hence $\rank(A) = n$. Because row operations do not change the rank, we have $\rank(A') = \rank(A)$, which implies that $\rank(A') = n$. The only $n \times n$ row-reduced echelon matrix with rank $n$ is the identity matrix, meaning $A' = I_n$. 

Since each step in the Gaussian elimination algorithm corresponds to an elementary operation, we can deduce from Lemma \ref{lemma:row_operations_matrix_multiplication} that there exist elementary matrices $T_1, \dots, T_k$ (of types I, II, and III) such that:
\begin{equation*}
    T_k \cdot T_{k-1} \dots T_1 \cdot A = I_n
\end{equation*}
By left-multiplying successively with the inverses of these elementary matrices, we isolate $A$ and deduce that $A = T_1^{-1} \dots T_k^{-1}$. Taking the inverse of $A$ yields $A^{-1} = T_k \dots T_1$.
\end{proof}

\begin{theorem}[Matrix Inversion Algorithm]
% thm:17.e.4
\label{thm:matrix_inversion_algorithm}
Let $A \in \M_{n \times n}(K)$. Construct the augmented matrix by placing $A$ alongside $I_n$, denoted as $(A \mid I_n)$. Then, $A$ is invertible if and only if $(A \mid I_n)$ can be brought, via a sequence of elementary row operations, to the form $(I_n \mid B)$ for some matrix $B \in \M_{n \times n}(K)$. Furthermore, if this is the case, we have $A^{-1} = B$.
\end{theorem}

Before proving this theorem, we require an intermediate technical result.

\begin{lemma}[Block Matrix Multiplication]
% lemma:17.e.5
\label{lemma:block_matrix_multiplication}
Let $B, C, D \in \M_{n \times n}(K)$. Then:
\begin{equation*}
    B \cdot \left(C \mid D\right) = \left(B \cdot C \mid B \cdot D\right)
\end{equation*}
\end{lemma}

\begin{proof}
This follows directly from the definition of matrix multiplication. Check the explicit details as an exercise!
\end{proof}

\begin{proof}[Proof of Theorem \ref{thm:matrix_inversion_algorithm}]
Assume that $A$ is invertible. By Theorem \ref{thm:elementary_matrix_decomposition}, there exist elementary matrices $T_1, \dots, T_k$ such that $T_k \dots T_1 \cdot A = I_n$. Using Lemma \ref{lemma:block_matrix_multiplication}, we can multiply the augmented matrix $(A \mid I_n)$ from the left by this chain of elementary matrices:
\begin{equation*}
    T_k \dots T_1 \cdot \left(A \mid I_n\right) = \left(T_k \dots T_1 \cdot A \mid T_k \dots T_1 \cdot I_n\right) = \left(I_n \mid T_k \dots T_1\right)
\end{equation*}
Let us define $B := T_k \dots T_1$. We can clearly see that:
\begin{equation*}
    B \cdot A = (T_k \dots T_1) \cdot A = I_n
\end{equation*}
and using the decomposition of $A$ provided by Theorem \ref{thm:elementary_matrix_decomposition},
\begin{equation*}
    A \cdot B = (T_1^{-1} \dots T_k^{-1}) \cdot (T_k \dots T_1) = I_n
\end{equation*}
This explicitly implies that $A^{-1} = B$.

Conversely, assume that $(A \mid I_n)$ can be brought to $(I_n \mid B)$ by a sequence of elementary row operations, for some matrix $B \in M_{n \times n}(K)$. This directly implies that the matrix $A$ itself is row-equivalent to $I_n$. Consequently, $A$ must have rank $n$, which inherently means that $A$ is invertible. This concludes the proof.
\end{proof}

\begin{example}
Let us find the inverse of the matrix $A = \begin{pmatrix} 1 & 2 \\ 3 & 4 \end{pmatrix} \in \M_{2 \times 2}(\mathbb{Q})$. We place $A$ and $I_2$ together in an augmented matrix and apply Gaussian elimination:
\begin{equation*}
    \left( \begin{array}{cc|cc} 1 & 2 & 1 & 0 \\ 3 & 4 & 0 & 1 \end{array} \right) \xrightarrow{R_2 - 3R_1 \longrightarrow R_2} \left( \begin{array}{cc|cc} 1 & 2 & 1 & 0 \\ 0 & -2 & -3 & 1 \end{array} \right)
\end{equation*}
Next, we scale the second row to create a pivot of $1$:
\begin{equation*}
    \xrightarrow{-\frac{1}{2}R_2 \longrightarrow R_2} \left( \begin{array}{cc|cc} 1 & 2 & 1 & 0 \\ 0 & 1 & \frac{3}{2} & -\frac{1}{2} \end{array} \right)
\end{equation*}
Finally, we eliminate the entry above the second pivot:
\begin{equation*}
    \xrightarrow{R_1 - 2R_2 \longrightarrow R_1} \left( \begin{array}{cc|cc} 1 & 0 & -2 & 1 \\ 0 & 1 & \frac{3}{2} & -\frac{1}{2} \end{array} \right)
\end{equation*}
The left side is now the identity matrix $I_2$, which means the right side is our inverse matrix. We conclude that:
\begin{equation*}
    A^{-1} = \begin{pmatrix} -2 & 1 \\ \frac{3}{2} & -\frac{1}{2} \end{pmatrix}
\end{equation*}
\end{example}